% Deutsche Parallelfassung zu runtime/runtime_operators.tex

\chapter{Laufzeitoperatoren}
\label{chap:runtime-operators}
\index{Laufzeitoperatoren}
\index{Operator!Laufzeit}

Laufzeitoperatoren definieren, wie etablierte Engineering Objects bei der
Ausführung interpretiert werden. Sie verwenden vorgelagerte Begriffe aus
Grundlagen und Zustand und bilden zwischen Interpretationsebenen ab. Sie
definieren Alignment-Identität, pose2/pose3, metrische oder physische
Realisierung und Repräsentation nicht neu.

\begin{principle}[Operator-Schichtprinzip]
\label{princ:operator-layer-principle}
Innerhalb von AIM wird Laufzeitinterpretation durch ausdrückliche Operatoren
organisiert, die verschiedene Zustandsschichten verbinden.
\end{principle}

\section{Operatorübersicht}

Die folgenden Symbole bezeichnen Familien qualifizierter Operationen in
späteren Kapiteln:
\[
\mathcal R,\qquad\mathcal D,\qquad\mathcal O,\qquad\mathcal V.
\]
Sie bilden keine verpflichtende Zustandsleiter. Eine Verwendung kann Soll mit
einer qualifizierten Realisierung vergleichen, Eisenbahnzustand mit \(s(t)\)
komponieren, ein gewähltes Modell auswerten, ein fahrzeugrelatives Ergebnis
ausdrücken oder eine zweckspezifische Bewertung anwenden. Jede Operation nennt
eigene Eingaben, Anwendbarkeit und Provenienz.

\section{Realisierungsoperator}

\begin{definition}[Realisierungsoperator]
\label{def:runtime-realization-operator}
\index{Realisierung!Operator}
Der Realisierungsoperator bildet Zielzustände auf realisierte Zustände ab:
\[
\mathcal R:x_{\mathrm{target}}\mapsto x_{\mathrm{real}}.
\]
\end{definition}
Für pose3 gilt
\[
\mathcal R_{\mathrm{pose3}}:
\mathrm{pose3}_{\mathrm{target}}\mapsto\mathrm{pose3}_{\mathrm{real}}.
\]
\(\mathcal R\) ist eine erklärende Kurzform für eine angegebene
Realisierungsbeziehung; der Operator macht Ausführung, physischen Zustand,
Instandhaltungsgeschichte und Beobachtung nicht zu einem Vorgang. Ein
gemessener Zustand bleibt Beobachtung mit eigener Provenienz und Unsicherheit.

\section{Dynamikoperator}

\begin{definition}[Dynamikoperator]
\label{def:runtime-dynamics-operator-operator-chapter}
\index{Dynamikoperator}
Für ein identifiziertes Antwortmodell \(M\) bildet der Dynamikoperator
qualifizierten Eisenbahnzustand, Durchfahrtsgesetz und Modelleingaben auf ein
modellspezifisches Ergebnis ab:
\[
\mathcal D_M:
(\mathcal S_{\mathrm{rail}}(s(t)),s(t),Q_M)\mapsto x_M(t).
\]
\end{definition}
Mögliche Ausgaben sind reduzierte Indikatoren oder modellierte Antworten.
Ihre Bezeichnungen begründen weder Komfort, Sicherheit noch Zulässigkeit. Eine
realisierungsbewusste Verwendung führt die gewählte qualifizierte Realisierung
als Eingabe von \(\mathcal D_M\) zu; Modell- und Durchfahrtsanforderungen
bleiben erhalten.

\section{Fahrzeuginterpretationsoperator}

\begin{definition}[Fahrzeuginterpretationsoperator]
\label{def:runtime-vehicle-operator}
\index{Fahrzeugzustand!Interpretationsoperator}
Für ein identifiziertes Fahrzeugmodell \(M\) bildet der Operator
qualifizierten Eisenbahnzustand, Durchfahrt und Fahrzeugeingaben auf
fahrzeugrelative Ergebnisse ab:
\[
\mathcal V_M:
(\mathcal S_{\mathrm{rail}}(s(t)),s(t),Q_M)
\mapsto x_{\mathrm{vehicle},M}(t).
\]
\end{definition}
Die qualifizierte Eingabemenge hält Fahrzeugtyp, Beladung, Modelltreue,
Provenienz und Anwendbarkeit getrennt, statt sie in einem einheitlichen
Kontext zu verbergen.

\section{Betrieblicher Bewertungsoperator}

\begin{definition}[Betriebsoperator]
\label{def:runtime-operational-operator}
\index{Betriebszustand!Bewertungsoperator}
Der Betriebsoperator bildet dynamische Laufzeitzustände auf betriebliche
Bewertungszustände ab:
\[
\mathcal O:x_{\mathrm{op}}\mapsto x_{\mathrm{eval}}.
\]
\end{definition}
\(\mathcal O\) bewertet Zulässigkeit, Normenkonformität, Komfortindikatoren
und Optimierungsziele, ohne die verwendeten vorgelagerten Zustandsobjekte neu
zu definieren.

\section{Operatorkomposition}
\index{Operatorkomposition}

Operatoren dürfen nur komponiert werden, wenn Ausgabe und Eingabe typmäßig
passen und Realisierung, Durchfahrt, Modell, Zweck, Provenienz und
Anwendbarkeit erhalten bleiben. Es gibt keine universelle Komposition, die
aus pose3 allein Fahrzeugantwort, Bewertung oder Entscheidung erzeugt.

\section{Laufzeitabstraktion und Implementierung}

\begin{principle}[Laufzeitabstraktionsprinzip]
\label{princ:runtime-abstraction-principle}
Laufzeitoperatoren definieren zunächst Interpretationsrollen; detaillierte
Mechanik kann später durch spezialisierte Implementierungen ergänzt werden.
\end{principle}

\begin{definition}[Operatorrolle]
\label{def:runtime-operator-role}
Eine Operatorrolle spezifiziert die begriffliche Transformation unabhängig von
einer bestimmten numerischen Implementierung.
\end{definition}

Diese Trennung hält AIM strukturell klar und implementierungsflexibel und
bewahrt zugleich die Kompatibilität mit Modellen höherer Genauigkeit.

\section{Kanonische Interpretation}

\begin{proposition}
Innerhalb von AIM ist Laufzeitinterpretation operatorbasiert.
\end{proposition}
\begin{proposition}
\(\mathcal R\) bildet Zielzustände auf realisierte Zustände ab.
\end{proposition}
\begin{proposition}
\(\mathcal D\) bildet Eisenbahn-Laufzeitzustände auf betriebliche dynamische
Zustandstrajektorien ab.
\end{proposition}
\begin{proposition}
\(\mathcal V\) bildet Eisenbahn-Laufzeitzustände auf fahrzeugrelative
Interpretationszustände ab.
\end{proposition}
\begin{proposition}
\(\mathcal O\) bildet Betriebszustände auf Bewertungs- und
Zulässigkeitszustände ab.
\end{proposition}

\section{Verbindung zu AIM}
\begin{summary}
Laufzeitoperatoren verbinden etabliertes Ingenieurwissen strukturell mit der
Laufzeitinterpretation. \(\mathcal R\) repräsentiert Realisierung,
\(\mathcal D\) dynamische Laufzeitzustandserzeugung, \(\mathcal V\)
fahrzeugrelative Interpretation und \(\mathcal O\) betriebliche Bewertung.
Ihre Komposition erzeugt abgeleitete Laufzeit-Ingenieurzustände, ohne
vorgelagerte Begriffe neu zu definieren.
\end{summary}
